\documentclass[12pt,a4paper]{article}
\usepackage{ctex}
\usepackage[utf8]{inputenc}
\usepackage[T1]{fontenc}
\usepackage{amsmath,amssymb}
\usepackage{graphicx}
\usepackage{booktabs}
\usepackage{hyperref}
\usepackage[margin=2.5cm]{geometry}
\usepackage{setspace}
\onehalfspacing

\title{透明但不按质量定价：首个代币化碳池崩溃的信用级证据}
\author{Wei Cai\textsuperscript{1,*}, Adeline Wen\textsuperscript{1}}
\date{}

\begin{document}
\maketitle

\noindent\textsuperscript{1} 去中心化计算实验室，工程与技术学院，华盛顿大学塔科马分校，塔科马，美国

\noindent\textsuperscript{*}通讯作者。邮箱：weicaics@uw.edu

\begin{abstract}
碳信用本应为真实的气候效益融资，但被当作等价商品出售的信用，实际带来的气候效益可能相差很大。项目可能即使没有碳收入也会建设；其减排主张可能难以核验，也可能无法持久。强弱信用以同一价格交易时，低质量信用最有动机进入。通常的应对办法是透明：公开信用的信息，让市场参与者自行区分质量。我们在Base Carbon Tonne（BCT）池检验这一思路。它是首个拥有完整公开交易记录的大型代币化碳市场。BCT将约2200万吨信用放到区块链上，任何人都可查验其底层记录，但每一笔信用都换成同一种可互换代币。我们据公共账本重建了该池完整生命周期：全部1,187笔存入、35,432笔提取，以及进行这些操作的账户。我们将每个代币与注册处记录相连，并从六项可观察特征评估底层信用，包括碳收入是否很可能是项目成立的必要条件，以及气候效益能否可靠监测。该池在2021至2026年间崩溃，常被归咎于森林保护类（REDD+）信用；但REDD+只占存入量的4.2\%。旧的可再生能源信用占69.1\%，其中许多项目即使没有碳收入也很可能会建设。少数账户提走了大部分受市场青睐的信用，留下930万吨低价值信用仍未被赎回、滞留池中。对于同时存在于BCT和一个限制存入资格的姊妹池中的信用，受限制池的赎回率低得多；这一描述性对比与池设计影响信用命运的解释一致。该案例表明，公开信息未必足以保护一个统一定价的异质资产池：参与者仍可在价格未反映质量差异时利用这些差异。因而，《巴黎协定》下的碳市场改革以及代币化真实世界资产的设计，可能不仅需要披露，也需要按质量区分价格的机制。
\end{abstract}

\section*{引言}
企业和政府购买碳信用，是为了支持这样一种主张：自己尚未避免的一吨排放，已在别处由一吨避免、减少或移除的排放所抵消。一张信用应代表一吨二氧化碳；但标称单位相同的两张信用，实际带来的气候效益可能大不相同。项目是否原本就会建设？所声称的减排能否被测量和独立核验？若碳被储存，它能在大气之外留多久？这些差异共同决定了信用质量。

其中最关键的是额外性：只有当碳收入改变了原本会发生的事情时，信用才具有额外性。一座即使没有碳收入也会建成的风电场，可以发电，也可以拥有完备文件，但出售该信用并不会额外带来气候效益。这个反事实本身无法直接观察。正因如此，在自愿碳市场（VCM）中区分强弱信用长期十分困难；VCM是企业自愿购买信用、而非为满足强制排放交易制度而购买的市场。其年交易额在2021至2022年曾突破20亿美元，但近期研究所审查项目中的信用，只有不到16\%代表真实减排\textsuperscript{4,5}。

经济学预言了市场无法区分质量时会出现的问题。Akerlof的“柠檬问题”描述的是：买方无法观察质量时，低质量商品会挤出高质量商品\textsuperscript{1,2,3}。碳市场常被置于这一框架下讨论。对此的设想是透明：让信用的历史和质量相关信息可见，使参与者避开差的信用、奖赏更好的信用。

区块链代币化承诺将这一设想付诸实践。它把交易记录在公共账本上，而不是记录在私下的场外双边交易中。最大的早期实验是Toucan协议的Base Carbon Tonne（BCT）池。持有人可将经Verra核证的信用送入Toucan的桥；Verra是全球最大的碳信用注册处。信用在链上被表示后进入BCT，持有人得到一种可互换的BCT代币。每一笔存入，无论底层信用质量如何，都换取同一种代币；所有代币按同一价格交易。之后，持有人可交回BCT代币，并选择提取某一特定底层信用。因此，池让质量相关信息可见，却没有让价格随质量变化。其代币价格从2021年底每吨8美元以上跌至2023年年中的0.71美元，到2026年已实际归零。

这一设计提出了传统碳市场难以直接回答的问题：参与者能够查阅质量相关信息时，仅靠透明能否阻止混合质量信用池恶化？传统市场很少留下能逐笔追踪信用从进入到退出的公共数据。BCT却不同：其账本记录了每笔存入、提取和账户。因此，它罕见地让研究者能够依据完整的可观察历史检验一种以透明为基础的市场设计，而不只是从价格或项目样本推断其行为。

既有的BCT崩溃解释也从未用这段历史检验。评论者和学者将崩溃归咎于低质量森林保护类信用，即REDD+，一种付费给土地所有者以避免毁林的项目\textsuperscript{2,6,7,8,9}。这一说法与REDD+信用问题的大量证据相符\textsuperscript{4,10,11}，但它并未说明这类信用是否真的主导BCT。更重要的是，它没有回答设计问题：如果池的组成不同，公开信息本身能否保护该池？

我们通过重建链上BCT来回答两个问题。我们收集所有1,187笔存入和35,432笔已记录的提取，识别509个存入账户和28,897个提取账户，并将每个信用代币与其注册处项目和年份相连。对于有逐项评分的161个代币，我们评估六项与气候质量相关的特征：它是移除还是避免排放，是否具有额外性，监测、报告与核验是否可靠，所声称效益是否持久，信用年份，以及使用的注册处和方法学。之后，我们用公共转账记录回答三个相连的问题：什么进入了池，什么被选择性提取，以及相同项目和年份类别的信用在设有筛选的姊妹池中是否有不同结局。这个组合使我们不仅能确认低质量信用是否存在，也能观察市场架构如何塑造它们的流动。

结果推翻了流行叙事。REDD+按存入吨数只占4.2\%。可再生能源信用占69.1\%，主要是来自中国、印度、土耳其和巴西的并网水电；现有证据表明，其中许多项目即使没有碳收入也会建设。少数账户提走了大部分高需求信用（155万吨），而930万吨低价值信用仍然滞留。对于同时进入BCT和设有质量筛选的姊妹池的14个信用代币类别，BCT中相同类别被全额提取，而筛选池中仅提取28.5\%。这是一项两池描述性对比，并非因果证明，但与池设计而非信用类型塑造资产命运的解释一致。

我们把这一机制称为\textit{设计使能的逆向选择}。它不是经典的柠檬问题：质量相关信息可以查阅，尽管参与者未必以同一方式解读。关键不对称性内嵌于池的规则：进入时，任何合格信用都换取同一通用代币；退出时，代币持有人可以选择取走哪一种信用。这就是盲目的进入与有选择的退出。统一价格不奖励向池中放入更好的信用，而退出时的选择权又使更好的信用值得被提走。完整账本让我们能在信用层面观察这一机制，这是不透明的链下市场无法提供的。

影响不止于一个失败的池。《巴黎协定》第6.4条正在发展新的联合国信用机制，而其他真实世界资产的代币化规模已超过250亿美元。两者都不等同于BCT；但只要一种机制将不同质量的资产以单一价格汇集，这一设计选择就既可检验也影响重大：透明可以为参与者提供信息，却可能仍不足以保护市场，还需要按质量区分价格以及对进入和退出两端设置规则。

\section*{结果}
\subsection*{什么进入了BCT，以及REDD+叙事遗漏了什么}

我们从BCT公共交易账本中重建了全部1,187笔存入。它们覆盖168个Verra核证碳标准（VCS）项目；VCS是一项记录已签发碳信用的注册处计划。总量约2200万吨（见方法）。

可再生能源信用主导了池（图1）。它们来自并网水电、风电和太阳能项目，以中国和印度为主，另有土耳其和巴西；这些项目注册于2008至2013年，合计占存入吨数的69.1\%。它们采用清洁发展机制（CDM）的类别，CDM是《京都议定书》下运行于2004至2020年的信用计划。自愿碳市场诚信委员会（ICVCM）是行业治理机构，未将这些类别批准为其核心碳原则（CCP）标签。计划层面的评估\textsuperscript{12,13}、风电CDM项目的荟萃分析\textsuperscript{4}以及对印度风电抵消的估计\textsuperscript{5}均发现额外性很低：这些项目很可能即使没有碳收入也会建设。



被广泛归咎为BCT崩盘元凶的REDD+信用，实际仅占池的4.2\%，不足可再生能源份额的十六分之一。即使将REDD+与其他基于自然的类别合并，即ARR（造林、再造林和植被恢复）和IFM（改良森林管理），总计也仅达10\%，而可再生能源单独就占69.1\%。


\subsection*{BCT的组成与Verra总体签发显著不同}

BCT的可再生能源主导可能仅反映Verra总体签发的供给。为此，我们将MSCI\textsuperscript{21}和Ecosystem Marketplace\textsuperscript{24}公布的VCS签发份额作为宽泛的供给参照：可再生能源约占按量计VCS签发的37\%（中心估计；敏感性范围26--48\%）。这一比较衡量BCT组成相对于注册处签发的差异，并不能识别任何存入者为何选择桥接某张信用。

近九成存入账户（89.0\%）的组合按吨数以可再生能源为主。BCT的可再生能源份额为69.1\%，是37\%注册处参照的1.87倍（比率为1表示没有组成差异）。一个地址可多次存入，故我们将地址而非交易视为独立单位；地址层面的随机化检验给出$p < 0.0001$（补充方法）。即使采用最高的48\%参照份额，差异仍为1.44倍；我们也检查了这一结果不是仅由少数大账户造成。

相反，REDD+以VCS基准率的0.2$\times$被低度代表（BCT 4.2\% vs. VCS约17\%），直接反驳了REDD+信用淹没池的叙事。

\textbf{偏斜从何处产生。}Toucan的桥是在信用进入池前、将注册处信用表示到链上的机制。以已桥接信用为条件，几乎所有桥接吨数都进入了BCT：按代币类别计93.5\%，按吨数计99.6\%（补充方法）。可再生能源偏斜因此产生于桥接决定之前或之时，而非BCT后续的准入筛选。账本不能区分是存入者选择还是其他因素使桥接流中可再生能源占比偏高。


\subsection*{价格与流入组成共同变动}

我们接下来检验价格与流入组成是否共同变动。我们将以美元计价的BCT每日价格（1,493个观测值，2021年10月至2026年7月；见方法）与从存入计算的累积组成指标合并（$n$ = 331个重叠天数）。这一观察性比较无法确定是什么驱动了价格下跌。

\textbf{共同变动。}价格与累积的吨数加权平均质量评分共同变动（Pearson $r$ = 0.77，$p < 10^{-66}$；图2a）。池级比较中，我们也报告其反向量，即池质量缺陷（PQD = 1 $-$ 评分/100）；PQD越高表示平均质量越低，它与价格的相关为$-$0.77。累积序列共享强烈的时间结构，故我们也比较逐日变化。一阶差分设定中，存入中可再生能源份额每增加一个百分点，与当日BCT价格低0.018美元相关（$p < 0.001$；$\beta$ = $-$1.8，份额以0至1表示；见方法）。探索性的周频时间先后检验更常将价格变化置于组成变化之前。这与价格下跌后流入信用质量降低相一致，但不能识别反馈效应或因果关系（补充方法）。

\textbf{什么离开了池。}35,432个向外转账事件表明，信用并非均匀地被提取。工业气体信用被全部赎回，REDD+为99.8\%，改良森林管理（IFM）为93.0\%，造林、再造林与植被恢复（ARR）为91.3\%；可再生能源信用则仅3.7\%。这一模式与退出端存在选择相一致。在最初逐项评分的161个代币子集中，已赎回信用的吨数加权平均复合评分为38.7（量表为0至100），已存入信用为31.7。整个生命周期中，BCT净可再生能源份额由71\%升至76\%。

提取顺序同样反映不同信用类型间的差异。在类型层面的比较中，较早被赎回的类型具有更高的链下市场需求（$\rho$ = $-$0.74，$n$ = 7种类型；补充方法）。2022年5月加密货币崩盘前，已赎回信用的吨数加权平均评分为42.0；之后为32.4。这一时序与“仍保有链下价值的信用在通用池代币成本较低时被提取”一致，但账本本身不揭示赎回者的动机。

存入者和赎回者几乎是完全不同的群体：28,897个赎回账户中仅1.4\%同时出现在509个存入账户中。前10名赎回者占提取吨数的85.0\%。他们的交易以快速爆发的方式进行（中位间隔4.6秒），与程序化执行一致。

账本区分了两项常被混为一谈的事实：BCT的准入规则没有筛选已被桥接的信用，而提取则明显集中于特定类型。前者描述宽松的架构，后者描述有选择的退出；它并不能证明同一批人，甚至同一批实体，推动了两个过程。


\subsection*{账户取证：谁提取了什么}

接下来，我们追踪161个最初评分代币类别的提取地址。最大的提取地址集中于特定信用类型，其中多数未曾出现在存入地址中。

\begin{table}[h!]
\caption{按吨数排名前5的赎回者。}\label{tab:top5}
\begin{tabular}{lllll}
\toprule
账户 & 赎回吨数 & 主导类型 & 平均质量 & 同时为存入者? \\
\midrule
0x65a5... & 651,334 & 工业气体 & 30.9 & 否 \\
0x1b8e... & 335,556 & IFM & 40.0 & 是 \\
0xee4b... & 195,051 & REDD+ & 31.4 & 否 \\
0x4b3e... & 188,205 & ARR & 50.4 & 否 \\
0x8556... & 183,629 & ARR & 46.0 & 否 \\
\bottomrule
\end{tabular}
\end{table}

这五个地址提取了155万吨：大多数已赎回工业气体信用、37\%的REDD+和44\%的ARR。最大的地址提取651,334吨工业气体信用（总量的42\%）。它是智能合约，即自主运行的链上程序；它在提取的同一笔交易中退役每一笔信用（经链上验证；见数据可用性）。这一模式与聚合器将信用导向最终抵消使用相一致。其余四个为外部拥有地址，不能据此认定为个人。它们提取后的转账与持有，与寻求价值的提取相一致，但不足以证明动机。中点价格假设下，潜在价差约为740万美元（若含退役合约为1000万美元；范围30万至2200万美元；见方法）。二十个最大提取地址中有十五个从未存入信用，且各自集中于一个类型。

追踪20个最大提取者所赎回的每一个信用的直接去向，显示出三条路径：跨池转移至一个设有质量筛选的池（NCT，Toucan的自然类限定池；39.2\%）、立即退役（34.6\%）和二级市场出售（17.0\%；补充方法）。

在既存入又提取的399个地址（1.4\%重叠）中，质量交换分析未发现“存入较弱信用而提取较强信用”的系统模式（补充方法）。观察到的两端因此主要涉及不同地址，但未必涉及完全不同实体。为检查这一限制，我们追踪两侧各前20个地址的首个可观察资金来源，发现一个共同资金源、涉及三个地址的直接信用转账以及两个两端均活跃的地址。这些可观察联系占前20存入吨数的26\%和前20提取吨数的9\%（补充表5）；它们无法解决交易所或中介背后隐藏的关联。

\subsection*{质量图谱与质量门控}

BCT在161个逐项评分代币子集中的池质量缺陷（PQD）为0.679。PQD等于吨数加权平均复合评分除以100后取反，故数值越高表示平均质量越低。它接近图3质量谱系的低端，该谱系从直接空气碳捕获与封存（DACCS）的0.076延伸至并网可再生能源的0.759。置换检验未发现存入子集的质量低于从我们定义的合格代币总体随机抽取的证据（$z$ = $-$0.64，$p$ = 0.27）；它不说明这些信用为何有资格或为何被桥接。我们将复合评分转为B至AAA等级（见方法）；BBB是我们预先规定、要求额外性获核实的第一个等级。另一项门控模拟使用完整的345个代币存入流，其中184个分数由匹配项目特征推算；BBB底线使PQD由0.689降至0.506，但只放行7.2\%的存入吨数（图4）。


\subsection*{时间动态与稳健性}

存入质量在BCT的运营生命周期内下降（Spearman $\rho$ = $-$0.24，$n$ = 1,187）。这一下降完全是年份效应：较晚的存入取自年代越来越久远的年份（年份即信用的减排年度），而非取自同一年份内更差的信用（补充表2）。2022年5月的加密货币崩盘提供了对所提反馈循环有用的时序检验。冲击后存入质量下降3.4个点（Cohen's $d$ = 0.49，中等效应），存入量下降98\%（图2b），质量下降斜率则增至三倍。这一模式与“价格先跌、存入质量随后恶化”一致，但并非一个干净的因果实验（补充方法）。

\subsection*{池设计与资产命运：一个代币内的跨池对比}

我们识别出31个同时存入无筛选BCT和Toucan Nature Carbon Tonne（NCT）池的项目--年份代币类别；NCT仅接纳基于自然的信用。随后，我们描述性地比较其中在BCT被全额提取的14个类别与其在NCT的赎回情况。它们共150万吨，而在NCT中仅28.5\%被赎回。逐类别的平均差异为+73.9个百分点（95\% CI [51, 93]；图5），每个纳入类别和信用类型方向一致。

这不是池设计效应的因果估计。14个类别是否进入样本取决于BCT端结果，BCT端的100\%因而部分由样本定义内建，+73.9个百分点是上界。两个池也在暴露时机、流动性、资格和参与者上不同。共同的项目--年份类别固定了项目类别和年份，却未消除这些池级差异。我们在补充方法报告精确检验和敏感性分析，作为对已观察配对的描述，而非独立分配处理的证据。

\textbf{筛选在机制上改变什么。}我们还在链上恢复了Toucan和第二个桥协议C3运营的五个池的信用身份，并用同一框架评分。平均评分随筛选严格程度上升：无筛选池约29--31，基于自然的筛选池约39--40，CHAR这一类别许可名单池为77.9（补充表4）。这不是筛选导致更高质量的独立检验。筛选由信用类型定义，而评分也主要由类型驱动，因此筛选会通过改变准入类型而机械地改变平均评分。筛选池与无筛选池中的基于自然信用在类型内评分几乎相同（差异不超过0.32分；补充方法）。该比较说明设计选择，但C3池没有外部验证，也无法识别同一类型内的效果。


\subsection*{赎回模式由信用类型描述得比复合等级更好}

捕获选择性赎回主导轴的，是信用类型本身，而非我们的质量框架。一个仅基于类型的预测规则优于质量等级规则；且在116个可再生能源代币内，年份和等级都无法预测赎回（补充表3）。因此，本文的核心发现，即组成、基准率过度选择、赎回端提取与代币内比较，建立在交易数据和Verra信用类型标签之上，而非评分框架之上（补充表6）。

结果是一大批滞留信用：在存入的960万吨B级信用中，有930万吨至今未被赎回。评分分析共覆盖1520万吨，这使BCT成为有记录以来最大的滞留碳信用池之一。

\section*{讨论}
在BCT中，公开的信用信息与交易记录未能阻止质量恶化。这不表示所有参与者都使用同一套质量衡量：注册处和账本公开的是来源与项目属性，而本文质量评估还结合了这些属性和已发表证据。该案例显示的是，当合约将异质信用视作可互换资产时，仅靠披露为何可能失效。

相关的经济直觉同时类似于“柠檬问题”和最便宜可交割激励。Akerlof的柠檬模型关心隐藏质量\textsuperscript{1}；BCT不同，因为用于评估信用的大量信息是公开的。我们所称的“设计使能的逆向选择”来自合约规则。账本重建提供了与Manshadi等人\textsuperscript{2}理论模型一致的描述性证据，但并未把两池对比变成因果识别。

所提出的机制很简单：若不同质量的资产以同一价格进入一个池，而其中一些在池外保有更高价值，持有人可能更愿意提走那些资产，较少受青睐的资产便会累积。披露可以揭示价差，却不消除价差。传统金融把相近激励记录为机构住房抵押贷款支持证券市场中的“最便宜可交割”定价\textsuperscript{25,26,27,28}。BCT额外提供了可观察的信用级进入与退出记录，使这一机制可以直接研究，而非仅由总价格推断。

将BCT的崩盘错误归因于REDD+信用\textsuperscript{7,8,9}，说明了结构性质量失败如何能隐藏于众目睽睽之下。CDM时代的可再生能源信用占BCT吨数的69.1\%。它们的额外性近乎为零，因为底层项目本身在经济上已具竞争力。与REDD+失败不同（后者围绕基线，即若没有碳信用项目本会出现的假设排放水平，以及毁林，产生可见争议），可再生能源信用的额外性失败在结构上是不可见的：大坝和风电场确实存在，电力也是真实的。聚焦于高知名度失败模式的质量保证框架，可能会错过主导池组成的那些类别。

超过900万吨名义上有效的抵消仍未被赎回（评分子集内非赎回率97.6\%）。在假定额外性范围、碳社会成本和注册处平均组成反事实的条件下，蒙特卡洛计算给出气候价值缺口中位数1.46亿美元（90\%模拟区间6800万至2.52亿美元；仅反映参数不确定性）。这是敏感性福利估计，而非观察到的损失。另一项基于价格假设的计算给出约1000万美元潜在私人价差（包括退役合约）；它不是已观察利润，也不能与福利估计相加。最大的提取地址退役了信用而非获取交易利润。上述数字分别描述池的组成、可能价差和流通结果。

这一机制可能延伸至碳以外的场景，但这一预测需要进一步检验。作为描述性比较，我们分析了NFTX的六个统一定价NFT金库，即独特数字收藏品池。存入物品的销售价值百分位在六个金库均低于所在系列中位数（中位百分位0.25--0.40，0.5表示无差异；补充方法）。没有金库显示相应的退出模式。这些数据表明另一场景中进入端的组成差异；它们不复制BCT的双边际机制，也不证明共同原因。

这些发现最直接关系到质量不同却以共同价格清算的自愿信用池与注册处捆绑\textsuperscript{14,15}。第6.4条是注册处信用机制，而非统一定价池；它继承的是设计原则，不是直接预测：只要质量未被定价，选择就可能出现。仅靠门控也未必足够：即使最低质量底线（复合评分$\geq$29）也会排除BCT 58\%的吨数。

该案例提出三项可供实现和检验的设计假设：按质量区分价格而非统一定价；同时考虑进入和退出规则；随组成改变复核质量底线。这些并非BCT单一案例已识别的效果。它们与ICVCM核心碳原则框架\textsuperscript{16}一致，也关系到新兴质量差异化系统，如新加坡将独立评级纳入资格评估\textsuperscript{17}以及区块链可信度指数\textsuperscript{18}。我们提供开源参考实现：链上注册表与门控、链下监测，单次门控检查约需19k gas\textsuperscript{19}（补充方法）。它是原型，不是部署后必能防止市场失灵的证据。

同样的数据说明了一种前瞻诊断，但尚未验证一项早期预警规则。累积柠檬指数是滚动的存入加权PQD，范围为0--1，越高越差。它在启动时为0.71。若采用说明性的0.50阈值，价格仍接近5.66美元启动水平时，指数已超过阈值，约早于价格跌破一半七个月。一个不依赖评分框架的替代指标，即累积可再生能源份额，也自首周起持续高于其阈值（补充方法）。因此信号是BCT的初始组成，而非之后的趋势。它启发一项可前瞻检验的协议：启动时审计组成，对存入和提取采用质量规则，并随池演变重新审计。

有几项局限值得说明。我们的质量评分由作者自行推导（商业评级仅覆盖极少数BCT项目），但核心发现依赖交易数据和Verra信用类型分类，而非评分框架。该框架本身已对照CCP标签（$d$ = 1.8；图6）、商业评级（$\rho$ = +0.901，$n$ = 27；这是一项非盲的一致性比较，因为扩展比较集的评分是在已知机构评级的情况下给出的；补充图1）和一个大语言模型评审团（Fleiss' $\kappa$ = 0.6，中等一致性；补充方法）进行了校准。两个跨池比较都是描述性的，而非因果识别。代币内对比（+73.9个百分点差距，$\Gamma$ = 3.05界）和池级双重差分都仅建立在两个池之上；它们的统计量量化的是每一配对的可变性，而非那些仍与筛选标签共线的池级混杂因素，即暴露时机、流动性和资格（见方法）。我们将这一对比解读为对池设计效应的强烈提示，并注意到占主导的可再生能源在结构上被排除在这一基于自然的比较之外。不过有一个潜在混杂可以排除：桥级分解（按吨数计99.6\%的传递率）表明进入端选择是架构性的而非策略性的。

政策含义可以证伪：当不同质量信用被自愿以单一价格汇集、且持有人可在退出时指定具体单位时，该机制预言市场存在脆弱性。ICVCM和第6.4条的透明度改革可以改善信息，但BCT表明，在这些条件下信息本身未必足够。第6.4条目前不是统一定价池；若它或其他机制将异质信用以单一价格汇集，可将这一预测与质量差异化定价或不同赎回规则的设计作比较检验。

代币化真实世界资产行业现已超过250亿美元，并正在跨商品、债券和应收账款复制基于池的架构。BCT是第一个以可审计交易数据完成完整质量崩溃周期的池。它的教训很直接：对于汇集碳市场而言，质量差异化的机制设计不是一项可选的改进，而是市场完整性的结构性前提。

\section*{方法}
（说明：英文投稿版按ERL快报体例仅保留方法摘要，以下完整细节在投稿版中位于补充材料；中文稿为阅读完整性保留全文。）

我们通过Polygon RPC端点采集了Polygon区块链上Toucan协议BCT池合约的所有存入交易。我们识别出1,187笔存入交易，时间跨度为BCT自2021年10月启动至2022年底新存入实际停止，映射至168个独立的VCS项目标识符（345个独立的TCO2代币地址），约2200万吨代币化碳信用。对于每笔存入，我们记录了：交易哈希、区块号和时间戳、存入者地址、TCO2代币合约地址、Verra VCS项目标识符、年份以及吨数。项目标识符与Verra注册处进行了交叉参照，以获取方法学类别、来源国家和CCP资格状态。

\textbf{项目分类。}每个项目根据Verra注册处条目被归入十一个方法学类别之一：可再生能源（并网风电、水电、太阳能；CDM方法学AMS-I.D、ACM0002）、化石燃料替代、废物管理和甲烷捕获、REDD+（VM0007、VM0009、VM0015）、造林/再造林（ARR）、工业气体销毁、改良森林管理（IFM）、能效提升、工业过程、农业和清洁炉灶。池级组成统计按吨数加权计算。

\textbf{NCT数据。}我们在相同的Polygon区块链上、相同的区块范围内采集了Toucan NCT池的全部708个存入事件。NCT将存入限制为AFOLU（农业、林业和其他土地利用）信用且年份$\geq$2012，作为对项目类型的一种最低质量筛选。NCT与BCT共享相同的协议基础设施和区块链底层。

\textbf{价格数据。}以USDC（一种与美元挂钩的稳定币）计价的每日BCT价格通过DeFi Llama API获取，数据来源为Polygon SushiSwap BCT-USDC交易对（1,493个日度观测值，2021年10月至2026年7月）。

\subsection*{质量评分框架}

我们开发了一个复合质量评分框架，评估六个活跃维度（选择这些维度，是为了覆盖信用完整性的主要决定因素；这些决定因素由牛津抵消原则和ICVCM CCP评估框架这两个最常被引用的规范性体系加以形式化），每个维度在0--100分的量表上评分：移除类型层级（权重0.250）、额外性（0.200）、MRV（监测、报告和核查）等级（0.200）、持久性（0.175）、年份（0.100），以及注册处与方法学质量（0.075）。初始权重来源于牛津抵消原则\textsuperscript{20}和ICVCM CCP评估框架\textsuperscript{16}。蒙特卡洛扰动（10,000个Dirichlet抽样权重向量）确认了93.7\%的等级稳定性。复合评分计算为$C = \sum_{i} w_i \times s_i$，并映射至字母等级：AAA（$\geq$90）、AA（$\geq$75）、A（$\geq$60）、BBB（$\geq$45）、BB（$\geq$30）、B（<30）。七个否决条件无论复合评分如何均设定最高等级上限。共同效益维度被赋予零权重，以避免通过共同效益叙事放大逆向选择\textsuperscript{6}。池质量缺陷（PQD = $1 - \bar{C}/100$）量化了每个池的吨数加权质量退化程度。完整的评分细节、分布式评分模型、扩展评分程序和敏感性分析见补充方法。

\subsection*{基准率比较}

BCT的可再生能源份额与Verra VCS基准率（中心估计37\%，敏感性范围26--48\%，来源为MSCI\textsuperscript{21}和Ecosystem Marketplace报告）进行了比较。我们将选择系数定义为比率$S = f_{\text{BCT,renewable}} / f_{\text{VCS,renewable}}$（为此目的而计算的一个量，而非借用的具名统计量），以账户聚类推断为主检验（置换、HHI调整和Bootstrap检验），朴素的精确二项检验仅作参照报告（见补充方法）。

\subsection*{代币内配对设计（描述性跨池对比）}

代币内池设计对比是根据同时存入BCT（无筛选）和NCT（设有质量筛选）的14个TCO2代币计算的。由于同一个物理信用同时出现在两个池中，质量、年份、注册处和项目类型在构造上保持不变，一个配对内唯一的差异便是池的架构；每个信用都是其自身的对照。两个池中每个代币的存入和赎回吨数直接从原始ERC-20 Transfer日志（存入 = 钱包$\rightarrow$池；赎回 = 池$\rightarrow$钱包）中重建，并按代币缓存，将已发布的总体赎回率（100.0\% / 28.5\%）复现至0.01个百分点以内，确认NCT的逐代币赎回率是测量得出的而非推算的。将代币作为分析单位（$n$ = 14），我们用三个精确程序（配对符号检验、对逐代币符号翻转的完全枚举置换检验，以及Wilcoxon符号秩检验）检验逐代币赎回率之差（BCT $-$ NCT），并用对14个逐代币差异的百分位Bootstrap（20,000次重采样）量化跨代币平均差距。精确检验的$p$值反映了差距方向的一致性，且无论是否剔除那个接近持平的代币（其BCT赎回量距全额赎回仅差约1吨）都保持不变。我们还拟合了一个吨级Beta--Binomial模型（存入吨数作为二项试验，均匀Beta(1,1)先验），但它假设各吨之间相互独立，而赎回实际上是分批集中发生的；它低估了不确定性，仅作描述性报告，推断则基于代币级Bootstrap和精确检验。估计对象是池设计对双重合格信用（同时被两个池接纳的信用）的效应，这是一个质量较高的基于自然的子集，而非完整的无筛选池。对未观测混杂的稳健性用Rosenbaum式敏感性界（符号检验结果失去显著性时的几率比$\Gamma$）加以量化。我们还拟合了一个带代币随机截距的混合效应logit模型，但它表现出准完全分离（BCT赎回率约100\%），仅作为方向性佐证报告。选择通道（哪些信用到达了两个池）由敏感性界和一个有向无环图加以界定（补充图3）。

\subsection*{池级双重差分（仅为描述性）}

我们还估计了一个汇集全部1,895个事件（1,187个BCT + 708个NCT）的存入级DiD：

$$\text{quality}_i = \beta_0 + \beta_1 \cdot \text{BCT}_i + \beta_2 \cdot \text{PostShock}_i + \beta_3 \cdot (\text{BCT} \times \text{PostShock})_i + \varepsilon_i$$

其中$\beta_3$捕获了BCT的质量在2022年5月崩盘后是否与NCT\textit{不同}地变化，标准误在TCO2代币层面聚类。我们仅将其报告为描述性佐证，而非因果估计：BCT和NCT共享其88.6\%的代币，违反了其推断所假设的单位独立性。在代币聚类Bootstrap下，标准误膨胀4.02倍，且池间秩差异与零无法区分（聚类稳健$p$ = 0.24，相对于朴素置换$p < 0.0001$）。因此，最强的基于设计的证据来自上文的代币内对比，我们将其解读为提示性而非识别性的。

\subsection*{多池泛化}

为评估该机制是否能泛化到单一池之外，我们用对BCT所用的同一方法学原型程序，对另外两个桥接工具（C3和Moss MCO2）的存入进行了评分。对于BCT和NCT，我们报告了同一运营方内部“无许可池与设有筛选池”对比下的时间-质量斜率（复合评分对存入顺序的Spearman $\rho$：$-$0.439对$-$0.105）。对于C3，逐代币的Verra项目和年份编码在代币符号中，并通过\texttt{eth\_call}在链上恢复（26个代币中22个真实；18个项目中17个已分类，其中五个经公共注册处查询；92笔存入中75笔已评分）；C3的时间-质量斜率为$\rho$ = $-$0.329（$p$ = 0.004），与BCT方向一致。Moss MCO2是一个单一的可替代篮子代币，没有逐代币的信用身份，因此逐代币质量信号确实无从定义，逆向选择在结构上不适用。剩余的覆盖缺口被明确记录而非推算。

\subsection*{赎回分析}

我们从BCT池合约作为发送方的ERC-20 Transfer事件中识别赎回，在161个最初评分的代币（2200万总吨数中的1520万吨）上共产生35,432个事件。赎回率按信用类型计算，为赎回吨数/存入吨数。代币内跨池比较利用了同时存入BCT和NCT的14个代币，以检验相同的信用在不同池设计下是否经历了不同的命运。账户取证、去向追踪和利润量化方法详见补充方法。

\subsection*{价格-质量动态}

累积PQD和可再生能源份额与每日BCT价格合并（$n$ = 331个重叠天数）。采用带有Newey--West HAC标准误（10个滞后期）的一阶差分OLS检验同期共变动。在周频率（$n$ = 55）下使用AIC选择的VAR(2)模型进行双向格兰杰因果检验。我们承认小样本限制了统计效力；每日回归（$n$ = 330）提供了更稳健的检验。

\subsection*{跨域比较}

为检验选择模式是否能泛化到碳之外，我们在六个NFTX金库（以太坊）上运行了与BCT同构的管线：全部铸入/赎回事件与钱包级金库代币转账由链上日志重建，每个被存入或被赎回的物品被赋予其最近一次成交的价格百分位，该百分位在该次成交前后一个滚动窗口内的同系列成交中计算（漂移受控的排名，基于排名稳定性假设；一个严格变体将匹配限制在事件本身窗口内的成交）。进入边际选择在每个金库内检验（Wilcoxon，存入物品百分位对照系列中位数）并跨金库检验（符号检验）；退出边际提取用Mann--Whitney检验赎回与存入百分位。处理某一侧超过5\%事件的地址被标记为路由合约并从账户重叠统计中排除。Curve stETH/ETH流量分析保留为支持性证据（观测性；亦与经典银行挤兑解释一致；补充方法）。

\subsection*{福利缺口估计}

我们使用蒙特卡洛模拟（10,000次迭代）估计了BCT质量退化的福利成本，选择该方法是因为多个不确定的输入（额外性比率、碳的社会成本、反事实组成）必须同时传播。每次迭代抽样：（i）基于文献分布的各信用类型额外性比率\textsuperscript{4,12,22,23}（可再生能源：0--20\%；REDD+：25--75\%；IFM：30--70\%；ARR：40--80\%）；（ii）碳的社会成本，来自均匀分布（\$50--200/tCO$_2$，涵盖美国EPA和斯特恩报告所用的范围）；（iii）匹配VCS注册处基准率的反事实池组成。福利缺口计算为反事实池与BCT实际组成之间预期气候价值的差异。该估计为说明性的，取决于额外性假设和反事实设定。

\subsection*{统计推断}

所有$p$值均使用Benjamini--Hochberg FDR程序在$\alpha$ = 0.05水平上对以下10项主要检验进行了校正：（1）基准率选择系数（账户聚类；二项检验仅作参照）、（2）全样本时间相关性、（3）冲击前时间相关性、（4）代币内配对符号检验、（5）池级池间差异（描述性，聚类稳健）、（6）格兰杰因果（质量$\rightarrow$价格）、（7）格兰杰因果（价格$\rightarrow$质量）、（8）池内置换检验、（9）存入者集中度Mann--Whitney检验、（10）赎回差异（代币级Mann--Whitney；$n$ = 161）。主要结论辅以10,000次置换检验$p$值。存入级统计使用聚类稳健Bootstrap（按存入者账户聚类，10,000次迭代）。额外的推断细节（账户聚类检验、CCP校准、评判者间信度、与商业机构的秩相关）见补充方法。

\subsection*{数据可用性}

所有数据、评分标准、分析脚本和智能合约源代码均可在 \url{https://github.com/Adeline117/carbon-neutrality} 获取（版本标签\texttt{v1.0-erl-submission}锁定本文所引用的状态），采用MIT许可证。机器可读的评分标准为JSON格式。包含逐维度评分的318个信用方法学批量数据集已纳入。所有345/345个BCT代币均已评分（161个原始 + 184个推算）。BCT（1,187笔交易）和NCT（708笔交易）的链上存入数据均附有交易哈希以支持独立验证。前几名赎回者的外部拥有账户与智能合约状态通过对一个公共Polygon节点使用\texttt{eth\_getCode}进行了检查；结果记录在\texttt{redeemer\_contract\_check.json}中。实体级独立性审计（两侧各前20个账户的首笔资金解析）可通过\texttt{entity\_funding\_analysis.py}复现，逐钱包结果见\texttt{entity\_funding\_analysis.json}。所有模型对29个信用的LLM评审团输出已纳入。分析脚本为纯Python，仅依赖NumPy/SciPy；单个校验脚本（\texttt{tools/verify\_headline\_numbers.py}）可对照缓存的分析输出重新断言每一个头条数字。


\subsection*{代码可用性}

所有分析代码均可在 https://github.com/Adeline117/carbon-neutrality 获取，采用MIT许可证。

\subsection*{作者贡献}

A.W.设计了研究方案，采集并分析了所有数据，开发了质量评分框架，进行了所有统计分析，并撰写了手稿。W.C.监督了研究并提供了关键修订。

\subsection*{利益冲突声明}

作者声明不存在利益冲突。

\subsection*{致谢}

本研究未获得任何资助机构的专项资助。作者感谢华盛顿大学塔科马分校去中心化计算实验室提供的计算资源和讨论。

\begin{thebibliography}{21}
\bibitem{1} Akerlof, G. A. The market for ``lemons'': quality uncertainty and the market mechanism. \textit{Q. J. Econ.} \textbf{84}, 488--500 (1970).
\bibitem{2} Manshadi, V. H. et al. Offsetting carbon with lemons: adverse selection and certification in the voluntary carbon market. SSRN Working Paper (2025).
\bibitem{3} Battocletti, V. et al. The voluntary carbon market: market failures and policy implications. \textit{Colo. Law Rev.} \textbf{95}, 519--572 (2024).
\bibitem{4} Probst, B. et al. Systematic assessment of the achieved emission reductions of carbon crediting projects. \textit{Nat. Commun.} \textbf{15}, 9562 (2024).
\bibitem{5} Calel, R. et al. Do carbon offsets offset carbon? \textit{Am. Econ. J. Appl. Econ.} \textbf{17}, 1--40 (2025).
\bibitem{6} Berg, F. et al. The market for voluntary carbon offsets. SSRN Working Paper (2025).
\bibitem{7} Bassi, E. et al. Blockchain-based voluntary carbon market: strategic insights into network structure. \textit{Front. Blockchain} \textbf{8}, 1603695 (2025).
\bibitem{8} Ballesteros-Rodriguez, A., De-Lucio, J. \& Sicilia, M.-A. Tokenized carbon credits in voluntary carbon markets: the case of KlimaDAO. \textit{Front. Blockchain} \textbf{7}, 1474540 (2024).
\bibitem{9} Jirasek, M. KlimaDAO: a crypto answer to carbon markets. \textit{J. Organ. Des.} \textbf{12}, 271--283 (2023).
\bibitem{10} Trencher, G. et al. Demand for low-quality offsets by major companies undermines climate integrity of the voluntary carbon market. \textit{Nat. Commun.} \textbf{15}, 6863 (2024).
\bibitem{11} West, T. A. P. et al. Action needed to make carbon offsets from forest conservation work for climate change mitigation. \textit{Science} \textbf{381}, 873--877 (2023).
\bibitem{12} Cames, M. et al. How additional is the Clean Development Mechanism? Oko-Institut Report (2016).
\bibitem{13} Schneider, L. Assessing the additionality of CDM projects: practical experiences and lessons learned. \textit{Clim. Policy} \textbf{9}, 242--254 (2009).
\bibitem{14} World Bank. State and Trends of Carbon Pricing 2025. World Bank Group (2025).
\bibitem{15} Garcia, A. \& Sanford, L. On the potential for strategic behavior in jurisdictional REDD+. \textit{Proc. Natl Acad. Sci. USA} \textbf{123}, e2531612123 (2026).
\bibitem{16} Integrity Council for the Voluntary Carbon Market. Core Carbon Principles, Assessment Framework, and Assessment Procedure. ICVCM (2023).
\bibitem{17} Singapore National Environment Agency. Carbon rating panel: appointment of BeZero, Calyx Global, and Sylvera. NEA Regulatory Notice (2025).
\bibitem{18} Gao, H. O. \& Liu, X. A blockchain-based carbon registry platform for credible climate action in transportation. \textit{npj Clim. Action} \textbf{5}, 23 (2026).
\bibitem{19} Reference implementation: composable quality-rating registry and gated pools (Solidity/Foundry). \url{https://github.com/Adeline117/carbon-neutrality} (contracts/), MIT license.
\bibitem{20} Allen, M. et al. The Oxford Principles for Net Zero Aligned Carbon Offsetting. University of Oxford (2020).
\bibitem{21} MSCI. State of integrity in the global carbon-credit market. MSCI ESG Research (2024).
\bibitem{22} Badgley, G., Freeman, J., Hamman, J. J., Haya, B., Trugman, A. T., Anderegg, W. R. L. \& Cullenward, D. Systematic over-crediting in California's forest carbon offsets program. \textit{Glob. Change Biol.} \textbf{28}, 1433--1445 (2022).
\bibitem{23} Haya, B. K. et al. Comprehensive review of carbon quantification by improved forest management offset protocols. \textit{Front. For. Glob. Change} \textbf{6}, 958879 (2023).
\bibitem{24} Ecosystem Marketplace. State of the Voluntary Carbon Market 2023. Forest Trends (2023).
\bibitem{25} Vickery, J. I. \& Wright, J. TBA trading and liquidity in the agency MBS market. \textit{FRBNY Econ. Policy Rev.} \textbf{19}, 1--18 (2013).
\bibitem{26} Fusari, N., Li, W., Liu, H. \& Song, Z. Asset pricing with cohort-based trading in MBS markets. \textit{J. Finance} \textbf{77}, 3249--3287 (2022).
\bibitem{27} An, Y., Li, W. \& Song, Z. TBA trading and security issuance in the agency MBS market. \textit{Real Estate Econ.} \textbf{53}, 607--642 (2025).
\bibitem{28} Huh, Y. \& Kim, Y. S. Cheapest-to-deliver pricing, optimal MBS securitization, and welfare implications. \textit{J. Financ. Econ.} \textbf{150}, 68--93 (2023).
\end{thebibliography}

\end{document}
